\documentclass[11pt, a4paper, UTF8]{chinese-erj}

% --- 必须的宏包 ---
\usepackage{amsmath, amssymb, mathtools}

% --- 动态引入本篇稿件的参考文献库 ---
\addbibresource{AERub_blinddate.bib}

% --- 动态定义出版信息 ---
\header{吉米奈、天马行空：最优停止与“正缘”匹配}
\pubinfo{2026年第2卷第1期} 

\begin{document}

% --- 衔接上一篇，页码从第3页开始 ---
\setcounter{page}{3} 

% ==========================================
% 标题与作者设置 (去除了寻租和框架效应，更加纯粹硬核)
% ==========================================
\title{最优停止与“正缘”匹配：基于37\%法则的相亲市场致胜策略\thanks{吉米奈，谷歌大中华区人工智能研发中心，邮政编码：100000，电子信箱：gemini@google.com；天马行空（通讯作者），小红书大学经济学院，邮政编码：200000，电子信箱：skyhorse@rednote.edu。本文受自燃科学鸡精委重点项目（项目批准号：72636007）“人类迷惑行为的经济学解释”资助。作者感谢匿名审稿专家的宝贵建议，当然文责自负。}}

\author{吉米奈 \quad 天马行空}
\date{}

\maketitle

\begin{abstract}
在当代中国相亲市场中，单身青年普遍面临“何时停止相亲”的动态决策难题。本文摒弃了传统的相亲焦虑叙事，将概率论中的经典“秘书问题”（Secretary Problem）转化为指导青年脱单的微观经济学致胜策略。研究推导表明，为了将遇到“全局最优伴侣（正缘）”的概率最大化，相亲者应科学运用“37\%法则”：将前37\%的相亲对象作为建立市场基准线的评估样本，随后一旦遇到超越该基准线的候选人便果断锁定。数学证明，即便在候选人基数趋于无穷大的极端不确定性下，该策略依然能将一击命中“正缘”的概率死死锁定在最高峰值36.8\%。本文为单身青年提供了一套打破迷茫、实现效用最大化的算法级相亲指南。

\keywords{相亲经济学 \quad 最优停止理论 \quad 37\%法则 \quad 匹配博弈 \quad 序贯决策}
\end{abstract}

% --- 正文 ---
\section{引言}

“我到底该选眼前这个还不错的人，还是继续等待下一个可能更好的人？”这是相亲市场中永恒的哈姆雷特之问。

长久以来，社会舆论总是用“挑花眼”、“剩女/剩男”等消极词汇来制造相亲焦虑。然而，相亲本质上是一个充满不确定性的序贯决策过程（Sequential Decision Process）。在信息不对称且无法“吃回头草”的规则下，盲目等待固然不可取，但草率决定同样会导致终身遗憾。

本文引入数学界著名的“秘书问题”\parencite{ferguson1989}，用严谨的最优停止理论（Optimal Stopping Theory）告诉广大青年：不要焦虑，数学早就为你计算出了相亲局里的“全局最优解”。只要掌握这个法则，你就能在茫茫人海中，以最大的概率精准狙击你的“正缘”。

\section{问题建模：相亲局的算法化}

假设在你的适婚年龄段内，你预期会遇到 $N$ 个潜在的相亲对象。为了求解这个最大化概率问题，我们设定如下模型条件：
\begin{enumerate}
    \item $N$ 个潜在伴侣按综合质量全序排列，且以完全随机的顺序与你见面。
    \item 每一轮相亲后，你必须立即做出排他性决策：“直接结婚”或“发好人卡并看下一个”。
    \item 你的目标极其宏大：不将就，不妥协，只选那个排名绝对第一的“最佳伴侣”（正缘）。
\end{enumerate}

面对如此苛刻的条件，我们需要寻找一个阈值 $k$：先无条件拒绝前 $k$ 个人，只观察不决策，以此建立你内心的“审美基准线”；从第 $k+1$ 个人开始，一旦出现比前面所有人都优秀的对象，立刻停止相亲，原地结婚。

\section{核心推导：寻找成功概率的巅峰}

\subsection{成功概率函数的构建}
在这个策略下，你选中最佳伴侣的概率 $P_{\text{成功}}(k)$，等于最佳伴侣出现在各个位置 $i$ 时的概率之和。

最佳伴侣出现在位置 $i$ 的先验概率为 $1/N$。要成功锁定位置 $i$（$i > k$）上的最佳伴侣，必须保证在前 $i-1$ 个人中，最优秀的那个人刚好出现在你的“观察期”（即前 $k$ 个人）内。否则，你会在到达位置 $i$ 之前，就被中间的某个次优项截胡。这一前置条件的概率为 $k/(i-1)$。

因此，命中正缘的总概率函数为：
\begin{equation}
    P_{\text{成功}}(k) = \sum_{i=k+1}^N \left( \frac{k}{i-1} \right) \cdot \frac{1}{N}
\end{equation}

\subsection{直击巅峰：$1/e$ 的数学奇迹}
为了找出让你成功率最高的那个“黄金分割点”，我们将离散的求和转化为连续的积分。令 $x = k/N$（观察期占总人数的比例）。当 $N$ 足够大时：
\begin{equation}
    P_{\text{成功}}(k) \approx \int_{x}^{1} \frac{x}{t} dt = -x \ln x
\end{equation}

对目标函数求一阶导数以寻找极大值：
\begin{equation}
    \frac{d}{dx}(-x \ln x) = -\ln x - 1
\end{equation}

令导数为零（$-\ln x - 1 = 0$），我们得到了宇宙级的相亲常数：
\begin{equation}
    x^* = e^{-1} \approx 0.367879
\end{equation}

这意味着，\textbf{最优的基准线截断点 $k^*$ 永远是总人数的 37\% 左右！}
更令人振奋的是，在这个最优策略下，你选中绝对“正缘”的最大化概率恰好也是：
\begin{equation}
    P_{\text{成功}}(k^*) \approx -e^{-1} \ln(e^{-1}) = e^{-1} \approx 36.8\%
\end{equation}

\section{结论：“36.8\%”不是焦虑，而是超能力}

在没有数学指导的盲目相亲中，如果你要在 100 个人里挑出唯一的那个“最对的人”，凭直觉的成功率只有微乎其微的 1\%。

然而，最优停止理论告诉我们：\textbf{只要你严格执行“前 37\% 只看不选，之后遇强则嫁/娶”的法则，无论你的相亲基数是 100 人还是 10000 人，你都能将找到极品真爱的概率，硬生生地拔高并锁定在惊人的 36.8\%！}

这不是一个让人沮丧的数字，这是一个数学奇迹。它赋予了单身青年在混沌的相亲市场中“逆天改命”的算法底气。

因此，我们的经济学建议是：不要再被长辈的催促打乱阵脚。评估一下你预期的总相亲次数，冷静地把前 37\% 的人当做你建立数据库的“市场调研”。一旦越过这个阈值，遇到那个打破你基准线的人，请毫不犹豫地果断出击。数学证明，那一刻，就是你这辈子离完美爱情最近的时刻。

% --- 输出参考文献 ---
\erjref
\printbibliography[heading=none]

\end{document}